\section{Discussion}\label{sec:discussion}

\paragraph{Summary.} Comparing classifiers on group-fairness criteria is expensive in labels because the relevant information sits in small cells: the positives and negatives of the minority group, and the instances on which the models disagree. PATTERN-DIF stratifies the pool by group and by the number of models predicting the positive class, allocates labels with a proxy-based generalised Neyman rule, and estimates the contrasts with a difference estimator that corrects the pool-wide proxy total by the labelled residuals. Across eleven data settings with nine models it reduces the labels needed for a given accuracy by a median factor of 5.8 at 250 labels, 6.3 at 500, 6.9 at 1,000 and 15.5 at 2,000 (5.5, 5.8, 6.0 and 9.4 without the degenerate law-school setting). The stratification contributes more than the estimator (a factor of 3.3 to 8.6 against 1.3 to 1.65), and the two are complementary. Stratifying by group alone gives no gain with the plain estimator. The design needs no label to set up and keeps the design-based guarantees of stratified sampling. Its benefit tends to be larger when the candidate models disagree less, which can be checked before any label is collected, although the rank correlation is only about $-0.6$ and the forecast is rough.

\paragraph{What the results do and do not show.} The efficiency figures are ratios of mean squared errors against simple random sampling on pools with known labels, so they give the number of labels the design saves for the same accuracy on the same pool. Eleven data sets are not representative of all comparisons an analyst will face. The gain depends on the disagreement of the compared models: it is between 3.7 and 5.8 at 500 labels for the settings with the highest disagreement (8.6\% to 9.5\%), and larger where the models are nearly identical. The estimator helps because the average predicted probability of nine models is informative in these data (area under the ROC curve 0.76 to 1.00) and, outside EEG, only moderately miscalibrated (10-bin calibration error 0.02 to 0.07; 0.20 in EEG). An uninformative or inverted proxy removes the advantage of the difference estimator over the plain one (Section~\ref{sec:ablations}), and the theory shows that the difference estimator loses whenever the within-stratum regression slope of the true influence on its proxy falls below one half (Section~\ref{sec:theory}). The allocation is insensitive to the proxy.

\paragraph{Relation to earlier work.} Sampling where models disagree is the central idea of active comparison of prediction models \citep{sawade2012active}, and verification designs that oversample discordant results have compared diagnostic tests in medicine for decades. Model-assisted estimation with predictions, including stratified versions, is the subject of prediction-powered inference \citep{angelopoulos2023ppi,fisch2024stratppi}. A stratified prediction-powered estimator that stratifies by the proxy score does not beat PATTERN-DIF in the paired median ratio of mean squared errors in any of the 35 cells of the main study (0.5 to 0.6 by budget, Section~\ref{sec:empirical}); its median efficiency against SRS exceeds that of PATTERN-DIF in one cell, Bank at 2,000 labels. Three features are specific to the present setting: the estimands are ratios whose numerators depend on disagreements and whose denominators depend on group prevalence everywhere; many contrasts across several models must be estimated from one sample; and the group is known before labelling. The stratification by the count of positive predictions, the allocation rule and the difference estimator together handle all three.

\paragraph{Limitations.} First, the main experiments use one split (three more for six settings), one set of nine models and pools of 1,841 to 19,537 instances per setting; larger and more diverse collections of models would change the disagreement structure and with it the gain. Second, in six of the eleven settings the group is arbitrary, so these settings test the statistics of the design and carry no fairness meaning. Third, the intervals are asymptotic. Mean coverage is 0.94 to 0.95, but at 250 labels the worst-covered contrast, in the law-school setting, has coverage 0.59. A variance floor for strata whose sampled labels are all equal raises the lowest value over settings from 0.59 to 0.82 (main study), and only a conservative version reaches the nominal level, with intervals 25\% wider in the median and 67\% wider on average (ablation study, Section~\ref{sec:ablations}). The floor was designed after seeing these failures, so its benefit is in-sample. Fourth, the primary estimands are signed gaps; the equalised-odds gap is a non-smooth function of them, and the design is a compromise for ranking models by it; the difference estimator fails for this purpose in the law-school setting (Section~\ref{sec:ablations}). Fifth, inference is about the fixed models on the given pool; variation from retraining is not addressed, and the two-phase interval for the deployment population was checked only on a synthetic population. Finally, labels are assumed correct; label noise that depends on the group would bias every design, including this one.

\paragraph{Extensions.} The design should extend to more than two groups (one stratum dimension per group level), to other ratio metrics such as positive predictive value or calibration within groups, and to sequential labelling, in which the allocation is updated as labels arrive by replacing the proxy moments with estimated ones. A stronger proxy, for instance a calibrated or learned model of the label given the predictions, could increase the gain of the difference estimator, which grows with the share of the label variance that the proxy explains (untested here). Combining the design with a retraining-aware bootstrap would address the variation due to model fitting.
